\section{Related Work}
\label{sec:related_work}

This section reviews two bodies of literature that underpin the present study: adaptive inertia weight strategies in PSO, with emphasis on fuzzy control approaches, and the influence of membership function design and linguistic granularity on fuzzy system performance.

\subsection{Adaptive Inertia Weight Control in PSO}
\label{subsec:rw_adaptive_w}

Since the introduction of the inertia weight by Shi and Eberhart~\cite{shi1998modified}, numerous strategies have been proposed to replace the original linear decreasing schedule with state-dependent adaptation. Chatterjee and Siarry~\cite{chatterjee2006nonlinear} introduced nonlinear decreasing functions that accelerate convergence in later iterations, while Nickabadi et al.~\cite{nickabadi2011novel} designed a feedback-based mechanism that adjusts~$w$ according to the swarm's success ratio, coupling the control signal to actual search progress rather than to a predefined temporal schedule.

Fuzzy logic provides a qualitative alternative by encoding adaptation knowledge as linguistic rules. Shi and Eberhart~\cite{shi2001fuzzy} first demonstrated that a Mamdani-type controller responding to swarm performance indicators can outperform static schedules, establishing a paradigm in which the inertia weight is governed by linguistic inference rather than an explicit function of the iteration count. Valdez et al.~\cite{valdez2017comparative} extended this approach with co-evolutionary strategies and showed consistent advantages over classical PSO across diverse benchmarks. Nobile et al.~\cite{nobile2018fuzzy} proposed a self-tuning mechanism that simultaneously adjusts multiple PSO parameters at the particle level, while Xia et al.~\cite{xia2022particle} introduced a multiple-input multiple-output fuzzy controller for adapting learning coefficients and elite selection. Olivas et al.~\cite{olivas2014comparative} showed that interval type-2 controllers can handle additional uncertainty from noisy landscapes, and Roy et al.~\cite{roy2024hyperparameter} recently proposed evolving fuzzy policies that adapt up to seven hyperparameters, representing the most comprehensive fuzzy control architecture for PSO reported to date.

A common feature across all these contributions is that the input membership functions---defining how diversity, fitness progress, or other indicators are fuzzified---are selected a priori and their geometric configuration is not examined as an experimental factor. Research efforts have focused on rule base design, output set selection, or the number of controlled parameters, treating the input partition as a fixed design element.

\subsection{Membership Function Design and Linguistic Granularity}
\label{subsec:rw_mf_design}

Evidence from the broader fuzzy systems literature suggests that the input partition is far from a neutral design choice. Olivas et al.~\cite{olivas2014comparative} demonstrated that different membership function shapes produce statistically significant performance differences even when the rule base remains unchanged, indicating that membership functions carry semantic information that modulates inference quality. Metaheuristic-based tuning of membership function parameters has confirmed this sensitivity across a range of applications~\cite{fierro2013optimization, nikolic2020bee}.

At the partition level, Cord\'on et al.~\cite{cordon2000uniform} established guidelines for obtaining uniform fuzzy partitions with adequate granularity, showing that properties such as overlap degree and label spacing are constrained by consistency requirements. Guillaume and Charnomordic~\cite{guillaume2004interpretable} proposed methods for generating interpretable partitions from data while preserving semantic coherence. These results establish that the geometric configuration of membership functions---support width, overlap, and boundary treatment---directly influences both the interpretability and the control behavior of the fuzzy system.

The number of linguistic labels used to partition the input space, referred to as linguistic granularity, constitutes a second design dimension. Castiello et al.~\cite{castiello2019interpretable} showed that adaptive granularity can improve discrimination in classification tasks without sacrificing interpretability. In optimization contexts, Huang et al.~\cite{huang2022linguistic} and Zhang et al.~\cite{zhang2022linguistic} demonstrated that the granularity of linguistic scales affects convergence behavior and solution quality in PSO-related decision-making problems.

Despite this evidence, the specific influence of input membership function geometry and linguistic granularity on fuzzy-controlled PSO remains unexplored. Existing fuzzy PSO studies adopt a single partition---typically three uniformly spaced triangular labels---without justifying this choice or evaluating alternatives. The present work addresses this gap through a factorial design in which four input set geometries, differing in overlap degree and boundary treatment, are crossed with two granularity levels (3 and 5~labels) under a constant output set and rule base. This design isolates the input partition as the sole experimental variable, enabling direct assessment of how fuzzification geometry affects optimization performance across problems with distinct landscape characteristics.